% Part: first-order-logic
% Chapter: tableaux
% Section: derivations

\documentclass[../../../include/open-logic-section]{subfiles}

\begin{document}

\iftag{FOL}
      {\olfileid{fol}{tab}{der}}
      {\olfileid{pl}{tab}{der}}

\olsection{\usetoken{P}{tableau}}

\begin{explain}
و مو ويل چې فرض څه دے، او د استنتاج قاعدې مو ورکړې۔ له دغو څخه
!!^{tableau}s په استقرايي ډول جوړېږي: هر !!{tableau} يا يوه يوازېنۍ
څانګه ده چې له يوه يا زياتو فرضونو جوړه وي، يا پر يوه څانګه د
استنتاج د کومې قاعدې په پلي کولو له بل !!a{tableau} څخه اخستل کېږي۔
\end{explain}

\begin{defn}[!!^{tableau}]
د فرضونو~$\sFmla{S_1}{!A_1}$، \dots،
$\sFmla{S_n}{!A_n}$ دپاره (چې هر~$S_i$ يا~$\True$ وي
يا~$\False$) !!^a{tableau} د~!!{signed formula}s يوه متناهي ونه
ده چې لاندې شرطونه پوره کوي:
\begin{enumerate}
\item د ونې تر ټولو پاسني $n$~!!{signed formula}s
  $\sFmla{S_i}{!A_i}$ دي، يو تر بل لاندې۔
\item په ونه کښې هر !!{signed formula} چې له فرضونو څخه نۀ وي، تر
  ځان پاس په څانګه کښې پر !!a{signed formula} د يوې استنتاجي قاعدې
  له سمې کارونې څخه راوتلے وي۔
\end{enumerate}
د~!!a{tableau} يوه څانګه هله او يوازې هله \emph{تړلې} ده چې
$\sFmla{\True}{!A}$ او~$\sFmla{\False}{!A}$ دواړه ولري؛ که نۀ وي
\emph{پرانيستې} ده۔ هغه !!^a{tableau} چې هره څانګه ئې تړلې وي، د
خپلو فرضونو د سټ \emph{تړلے !!{tableau}} دے۔ که !!a{tableau} تړلے
نۀ وي، يعنې لږ تر لږه يوه پرانيستې څانګه ولري، نو \emph{پرانيستے}
دے۔
\end{defn}

\begin{ex}
  د فرضونو هر سټ په يوازې ځان !!a{tableau} دے، خو په عام ډول به
  تړلے نۀ وي۔ (په څرګند ډول يوازې هله تړلے دے چې فرضونه له مخکښې
  د~!!{signed formula}s يوه جوړه
  $\sFmla{\True}{!A}$ او~$\sFmla{\False}{!A}$ ولري۔)

  له يوه !!a{tableau} (پرانيستي يا تړلي) څخه د کومې استنتاجي قاعدې
  په پلي کولو يو نوے، لوے !!{tableau} اخستلے شو۔ قاعده په هغه کښې پر
  !!a{signed formula}~$!A$ پلې کوو۔ دا قاعده به د هرې هغې څانګې په
  پاے کښې يو يا زيات !!{signed formula}s ورزيات کړي چې د~$!A$ هغه
  پېښه لري چې قاعده مو پرې پلې کړې ده۔

  مثلاً فرض~$\sFmla{\True}{!A \land \lnot !A}$ ونيسئ۔ لاندې هغه
  (پرانيستے) !!{tableau} دے چې يوازې له همدې فرض څخه جوړ دے:
  \begin{center}
    \begin{tableau}{}
      [\sFmla{\True}{\formula{A} \land \lnot \formula{A}}, just=\TAss]
    \end{tableau}{}
  \end{center}
  د~$\TRule{\True}{\land}$ قاعده پر فرض پلې کوو او ترې نوے
  !!{tableau} اخلو۔ دا قاعده راکوي چې تابلو ته دوه نوې کرښې،
  $\sFmla{\True}{!A}$ او~$\sFmla{\True}{\lnot !A}$، ورزياتې کړو:
  \begin{center}
    \begin{tableau}{}
      [\sFmla{\True}{\formula{A} \land \lnot \formula{A}}, just=\TAss
        [\sFmla{\True}{\formula{A}}, just={\TRule{\True}{\land}[1]},
          [\sFmla{\True}{\lnot \formula{A}}, just={\TRule{\True}{\land}[1]}
          ]
        ]
      ]
    \end{tableau}{}
  \end{center}
  کله چې !!{tableau}s ليکو، پلې شوې قاعدې په ښي لوري کښې ثبتوو
  (مثلاً $\TRule{\True}{\land} 1$ دا معنٰی لري چې پر دې کرښه
  !!{signed formula} د~$\TRule{\True}{\land}$ قاعدې د هغې کارونې
  پايله ده چې د~$1$ پر کرښه !!{signed formula} ته شوې ده)۔ دا نوے
  !!{tableau} اوس نور !!{signed formula}s لري، خو يوازې پر يوه
  ($\sFmla{\True}{\lnot !A}$) قاعده پلې کولے شو (په دې حالت کښې
  $\TRule{\True}{\lnot}$ قاعده)۔ له دې څخه لاندې تړلے
  !!{tableau} راځي:
  \begin{center}
    \begin{tableau}{}
      [\sFmla{\True}{\formula{A} \land \lnot \formula{A}}, just=\TAss
        [\sFmla{\True}{\formula{A}}, just={\TRule{\True}{\land}[1]}
          [\sFmla{\True}{\lnot \formula{A}}, just={\TRule{\True}{\land}[1]}
            [\sFmla{\False}{\formula{A}}, just={\TRule{\True}{\lnot}[3]}, close]
          ]
        ]
      ]
    \end{tableau}{}
  \end{center}
\end{ex}

\end{document}

